\documentclass[12pt,a4paper]{article}
\usepackage{amsmath, amssymb}
\usepackage[margin=2.5cm]{geometry}
\usepackage{enumitem}
\usepackage{xcolor}

\newcommand{\easy}{\textcolor{green!60!black}{$\bigstar$}}
\newcommand{\medium}{\textcolor{orange}{$\bigstar\bigstar$}}
\newcommand{\hard}{\textcolor{red}{$\bigstar\bigstar\bigstar$}}

\title{\textbf{Solutions: Additional Exercises 4}\\Taylor Approximations \& Uniform Convergence}
\author{Series and Integral Transforms}
\date{}

\begin{document}
\maketitle

% -------------------------------------------------------
\section*{Exercise 1 — Numerical Approximations via Taylor Series}

\textit{General strategy}: Use the Taylor/Maclaurin series, bound the remainder $R_n$ by the first omitted term (for alternating series, or by the Lagrange remainder otherwise), and find the smallest $n$ that makes $|R_n|<\varepsilon$.

\subsection*{(a) \easy\ $e^{0.1}$ to within $10^{-4}$}

$e^x = \sum_{k=0}^{n}\dfrac{x^k}{k!} + R_n(x)$ where $|R_n(x)|\leq \dfrac{e^{|x|}|x|^{n+1}}{(n+1)!}$.

For $x=0.1$, bound $e^{0.1}<3$:
\[
  |R_n|\leq \frac{3\cdot(0.1)^{n+1}}{(n+1)!}.
\]
For $n=3$: $|R_3|\leq \dfrac{3\cdot(0.1)^4}{4!}=\dfrac{3\times10^{-4}}{24}\approx 1.25\times10^{-5}<10^{-4}$.

\textbf{Use $n=3$ terms}:
\[
  e^{0.1}\approx 1+0.1+\frac{(0.1)^2}{2}+\frac{(0.1)^3}{6}
  = 1+0.1+0.005+0.000\overline{1}67
  \approx \mathbf{1.1052}.
\]

\subsection*{(b) \medium\ $\sin(0.2)$ to within $10^{-5}$}

$\sin x = \sum_{k=0}^{n}\dfrac{(-1)^k x^{2k+1}}{(2k+1)!}$ (alternating).  The error is bounded by the first omitted term:
\[
  |R_n| \leq \frac{(0.2)^{2n+3}}{(2n+3)!}.
\]
For $n=2$ (using terms $k=0,1,2$, i.e., $x, -x^3/6, x^5/120$):
\[
  |R_2|\leq\frac{(0.2)^7}{5040}\approx\frac{1.28\times10^{-5}}{5040}\approx2.5\times10^{-9}\ll10^{-5}.
\]
Even $n=1$ gives $|R_1|\leq(0.2)^5/120\approx 2.7\times10^{-6}<10^{-5}$.  \textbf{Use $n=1$ (2 non-zero terms)}:
\[
  \sin(0.2)\approx 0.2 - \frac{(0.2)^3}{6}
  = 0.2 - 0.001\overline{3}
  \approx \mathbf{0.19867}.
\]

\subsection*{(c) \medium\ $\ln(1.1)$ to within $10^{-3}$}

$\ln(1+x)=\sum_{k=1}^{\infty}\dfrac{(-1)^{k+1}x^k}{k}$ at $x=0.1$ (alternating, since $0<x\leq 1$).

Error bounded by first omitted term: $|R_N|\leq\dfrac{(0.1)^{N+1}}{N+1}$.

For $N=3$: $|R_3|\leq\dfrac{(0.1)^4}{4}=2.5\times10^{-5}<10^{-3}$.

\textbf{Use 3 terms}:
\[
  \ln(1.1)\approx 0.1-\frac{0.01}{2}+\frac{0.001}{3}
  = 0.1-0.005+0.000\overline{3}
  \approx \mathbf{0.09531}.
\]

\subsection*{(d) \medium\ $\sqrt[3]{28}$ to within $10^{-3}$}

Write $28=27\cdot(1+1/27)$, so $\sqrt[3]{28}=3\cdot(1+1/27)^{1/3}$.

Use $(1+u)^{1/3}=\sum_{k=0}^{\infty}\binom{1/3}{k}u^k$ at $u=1/27$:
\[
  (1+u)^{1/3}\approx 1+\frac{u}{3}-\frac{u^2}{9}+\cdots
\]
For $u=1/27$: $u/3=1/81\approx0.01235$; $u^2/9=1/6561\approx1.5\times10^{-4}$.

The second correction term is $\approx4.6\times10^{-4}$.  Using two correction terms:
\[
  \sqrt[3]{28}\approx 3\left(1+\frac{1}{81}-\frac{1}{6561}\right)
  = 3\times1.012195\ldots\approx\mathbf{3.0366}.
\]

\subsection*{(e) \medium\ $\cos(\pi/36)$ to within $10^{-8}$}

$\pi/36\approx0.08727$.  $\cos x=\sum_{k=0}^{\infty}\dfrac{(-1)^k x^{2k}}{(2k)!}$ (alternating); error $\leq$ first omitted term.

For $n=3$ (terms $k=0,1,2$):
\[
  |R_3|\leq\frac{(\pi/36)^6}{720}\approx\frac{(0.08727)^6}{720}\approx\frac{4.4\times10^{-7}}{720}\approx6\times10^{-10}<10^{-8}.
\]
\textbf{Use 3 terms}:
\[
  \cos(\pi/36)\approx 1-\frac{(\pi/36)^2}{2}+\frac{(\pi/36)^4}{24}
  \approx 1-0.003805+7.24\times10^{-6}
  \approx\mathbf{0.99620}.
\]

\subsection*{(f) \hard\ $\pi/4=\arctan(1)$ via series}

$\arctan x=\sum_{k=0}^{\infty}\dfrac{(-1)^k}{2k+1}x^{2k+1}$.  At $x=1$, this is the Leibniz series $1-1/3+1/5-\cdots$.

Error: $|R_N|\leq 1/(2N+3)<10^{-3}$ requires $2N+3>1000$, i.e., $N\geq499$. This is extremely slow.

\textit{Faster approach}: use $\pi/4=\arctan(1/2)+\arctan(1/3)$ (Machin-type identity).  Then each series converges much faster; only $\sim 3$–$4$ terms are needed to achieve $10^{-3}$ accuracy.

\subsection*{(g) \hard\ $\ln 2$ efficiently}

The direct series $\ln 2=\sum_{k=1}^{\infty}(-1)^{k+1}/k$ converges, but slowly ($|R_N|\leq 1/(N+1)$, need $N\geq999$).

\textit{Better}: use $\ln 2 = 2\operatorname{arctanh}(1/3)$ where $\operatorname{arctanh}(u)=\sum_{k=0}^{\infty}\dfrac{u^{2k+1}}{2k+1}$:
\[
  \ln 2 = 2\sum_{k=0}^{\infty}\frac{(1/3)^{2k+1}}{2k+1}.
\]
Error after $N$ terms: $|R_N|\leq \dfrac{2\cdot(1/3)^{2N+3}}{2N+3}$.

For $N=2$: $|R_2|\leq \dfrac{2\cdot(1/3)^7}{7}\approx 3.2\times10^{-4}<10^{-3}$.  Only 3 terms needed. $\checkmark$

\subsection*{(h) \hard\ $\displaystyle\int_0^{0.5}e^{-x^2}\,dx$ to within $10^{-4}$}

Use $e^{-x^2}=\sum_{k=0}^{\infty}\dfrac{(-1)^k x^{2k}}{k!}$ and integrate term by term:
\[
  \int_0^{0.5}e^{-x^2}\,dx = \sum_{k=0}^{\infty}\frac{(-1)^k (0.5)^{2k+1}}{k!(2k+1)}.
\]
This is an alternating series; error $\leq$ first omitted term.

For $N=3$ (terms $k=0,1,2,3$): first omitted term ($k=4$):
\[
  \frac{(0.5)^9}{4!\cdot 9} = \frac{(1/2)^9}{216} = \frac{1}{110592}\approx 9.0\times10^{-6}<10^{-4}.
\]
\textbf{Use 4 terms} ($k=0,1,2,3$):
\[
  \approx 0.5 - \frac{(0.5)^3}{3} - \frac{(0.5)^5}{2!\cdot5} + \frac{(0.5)^7}{3!\cdot7}
  \approx 0.5 - 0.04167 + 0.003125 - 0.000186\ldots
  \approx \mathbf{0.4612}.
\]

% -------------------------------------------------------
\section*{Exercise 2 — Uniform Convergence of Sequences}

\subsection*{(a) \easy\ $f_n(x)=x^n$ on $[0,1)$}

Pointwise limit: $f(x)=0$ for $x\in[0,1)$.

$\sup_{x\in[0,1)}|f_n(x)-0|=\sup_{x\in[0,1)}x^n = 1\not\to 0$.

\textbf{NOT uniformly convergent} on $[0,1)$.

\subsection*{(b) \medium\ $f_n(x)=x/n$ on $\mathbb{R}$}

Pointwise limit: $0$.  But $\sup_{x\in\mathbb{R}}|x/n|=+\infty$.  \textbf{NOT uniform on $\mathbb{R}$}.

On any bounded interval $[a,b]$: $\sup|x/n|\leq \max(|a|,|b|)/n\to 0$.  \textbf{Uniform on every bounded interval}.

\subsection*{(c) \easy\ $f_n(x)=\dfrac{1}{nx+1}$ on $[1,\infty)$}

For $x\geq 1$: $0\leq f_n(x)\leq\dfrac{1}{n+1}\to 0$.  So $\sup_{x\geq 1}|f_n(x)|\leq\dfrac{1}{n+1}\to 0$.

\textbf{Uniformly convergent} to $0$ on $[1,\infty)$.

\subsection*{(d) \medium\ $f_n(x)=\dfrac{nx}{1+n^2x^2}$ on $(0,\infty)$}

Pointwise: for fixed $x>0$, $f_n(x)=\dfrac{nx}{1+n^2x^2}\leq\dfrac{1}{2nx}\to 0$ (by AM-GM: $1+n^2x^2\geq 2nx$).

But at $x_n=1/n$: $f_n(1/n)=\dfrac{n\cdot(1/n)}{1+1}=\dfrac{1}{2}$.

So $\sup_{x>0}|f_n(x)|\geq 1/2\not\to 0$.  \textbf{NOT uniformly convergent}.

\subsection*{(e) \hard\ $\displaystyle\int_0^1 n^2 x\,e^{-nx}\,dx$}

Integrate by parts ($u=x$, $dv=n^2e^{-nx}dx$):
\[
  \int_0^1 n^2 x e^{-nx}\,dx
  = \left[-nxe^{-nx}\right]_0^1 + \int_0^1 ne^{-nx}\,dx
  = -ne^{-n} + \left[-e^{-nx}\right]_0^1
  = -ne^{-n} + 1 - e^{-n}
  = 1-(n+1)e^{-n}\to 1.
\]
The pointwise limit of $f_n(x)=n^2xe^{-nx}$ is $0$ on $(0,1]$, and $\int_0^1 0\,dx=0\neq 1$.

Since $\displaystyle\lim_{n\to\infty}\int_0^1 f_n\,dx \neq \int_0^1\lim_{n\to\infty}f_n\,dx$, the convergence is \textbf{NOT uniform} on $[0,1]$.

\subsection*{(f) \medium\ $f_n(x)=\dfrac{\sin(nx)}{\sqrt{n}}$}

$\sup_x|f_n(x)|=1/\sqrt{n}\to 0$.  \textbf{Uniformly convergent} to $0$ on $\mathbb{R}$.

However, $f_n'(x)=\sqrt{n}\cos(nx)$, which is unbounded.  The convergence of derivatives fails — a warning that uniform convergence does \textbf{not} imply convergence of derivatives.

\subsection*{(g) \medium\ $f_n(x)=x^n(1-x)$ on $[0,1]$}

Pointwise limit: $0$.  Find the maximum on $[0,1]$:
$f_n'(x)=x^{n-1}(n-(n+1)x)=0$ at $x^*=n/(n+1)$.
\[
  f_n(x^*) = \left(\frac{n}{n+1}\right)^n\cdot\frac{1}{n+1}.
\]
Since $\left(\dfrac{n}{n+1}\right)^n\to e^{-1}$:
\[
  \sup_{[0,1]}|f_n|\sim\frac{e^{-1}}{n}\to 0.
\]
\textbf{Uniformly convergent} to $0$ on $[0,1]$.

\subsection*{(h) \medium\ Proof: uniform convergence implies integral convergence}

\textit{Theorem}: If $f_n\to f$ uniformly on $[a,b]$, then $\int_a^b f_n\to\int_a^b f$.

\textit{Proof}: For any $\varepsilon>0$, there exists $N$ such that $\|f_n-f\|_\infty<\dfrac{\varepsilon}{b-a}$ for all $n>N$.  Then:
\[
  \left|\int_a^b f_n\,dx - \int_a^b f\,dx\right|
  = \left|\int_a^b(f_n-f)\,dx\right|
  \leq \int_a^b|f_n-f|\,dx
  \leq (b-a)\cdot\frac{\varepsilon}{b-a} = \varepsilon. \quad\square
\]

% -------------------------------------------------------
\section*{Exercise 3 — Uniform Convergence of Series}

\subsection*{(a) \easy\ $\displaystyle\sum\frac{\cos(nx)}{n^2}$ on $\mathbb{R}$}

$\left|\dfrac{\cos(nx)}{n^2}\right|\leq\dfrac{1}{n^2}:=M_n$ and $\sum M_n=\sum1/n^2<\infty$.

By the \textbf{Weierstrass $M$-Test}: \textbf{uniformly convergent on $\mathbb{R}$}.

\subsection*{(b) \medium\ $\displaystyle\sum x^n$ on $(-1,1)$ vs. $[-r,r]$, $0<r<1$}

On $[-r,r]$: $|x^n|\leq r^n=:M_n$ and $\sum r^n<\infty$.  \textbf{Uniform by $M$-Test}.

On $(-1,1)$: $\sup_{x\in(-1,1)}|x^n|=1$ and $\sum 1=\infty$.  The $M$-Test fails. In fact, the convergence is \textbf{not uniform} on $(-1,1)$ (the sum $1/(1-x)$ is unbounded near $x=1$).

\subsection*{(c) \medium\ $\displaystyle\sum\frac{(-1)^n x^{2n}}{n}$ on $[-1,1]$}

Use \textbf{Dirichlet's Test} for uniform convergence: the partial sums of $(-1)^n$ are bounded (by $1$), and $x^{2n}/n$ decreases to $0$ uniformly on $[-1,1]$ (since $|x^{2n}/n|\leq 1/n\to 0$ uniformly).  \textbf{Uniformly convergent on $[-1,1]$}.

\subsection*{(d) \medium\ $\displaystyle\sum\frac{x^n}{n^2+x^2}$}

On $[-r,r]$: $\dfrac{|x|^n}{n^2+x^2}\leq\dfrac{r^n}{n^2}=:M_n$ and $\sum r^n/n^2<\infty$.  \textbf{Uniform on $[-r,r]$} by $M$-Test.

On $[-1,1]$: at $x=1$, $a_n=1/(n^2+1)$ and $\sum 1/(n^2+1)$ converges.  A careful analysis shows uniform convergence holds on $[-1,1]$ as well.

\subsection*{(e) \medium\ $S(x)=\displaystyle\sum\frac{\sin(nx)}{n^3}$: continuity and differentiability}

$|\sin(nx)/n^3|\leq 1/n^3$; $\sum 1/n^3<\infty$.  By $M$-Test, the series converges \textbf{uniformly on $\mathbb{R}$}, so $S$ is \textbf{continuous}.

Differentiated series: $\sum\dfrac{n\cos(nx)}{n^3}=\sum\dfrac{\cos(nx)}{n^2}$.  By part (a), this converges uniformly.  By the \textbf{Term-by-Term Differentiation Theorem}:
\[
  S'(x) = \sum_{n=1}^{\infty}\frac{\cos(nx)}{n^2}.
\]

\subsection*{(f) \medium\ $\displaystyle\sum_{n=0}^{\infty}(-1)^n x^{2n}$ on $[-r,r]$, $r<1$}

$|(-1)^n x^{2n}|\leq r^{2n}$; $\sum r^{2n}=1/(1-r^2)<\infty$.  \textbf{Uniform by $M$-Test}.

The series sums to the geometric series $\sum(-x^2)^n=\dfrac{1}{1+x^2}$.

Differentiating term by term (justified by uniform convergence on $[-r,r]$):
\[
  \frac{d}{dx}\!\left(\frac{1}{1+x^2}\right) = -\frac{2x}{(1+x^2)^2}
  = \sum_{n=1}^{\infty}(-1)^n(2n)x^{2n-1}. \checkmark
\]

\subsection*{(g) \hard\ Power series $\sum a_n x^n$ on $[-r,r]$ vs. full $(-R,R)$}

\textit{Uniform on $[-r,r]$ for any $r<R$}: $|a_n x^n|\leq|a_n|r^n=:M_n$.  Since $r<R$, the series $\sum|a_n|r^n$ converges (by definition of radius of convergence and Abel's theorem).  \textbf{$M$-Test applies}: uniform convergence on $[-r,r]$. $\square$

\textit{Not necessarily uniform on $(-R,R)$}: Consider $\sum x^n$ with $R=1$.  The sum $1/(1-x)\to\infty$ as $x\to 1^-$, so the partial sums cannot converge uniformly on $(-1,1)$ (a uniformly convergent series of bounded functions has a bounded sum).

\bigskip
\hrule
\bigskip
\textit{Key takeaways}: For sequences, compute $\sup_x|f_n(x)-f(x)|$ and check if it tends to $0$; a moving maximum that stays bounded away from $0$ kills uniform convergence. For series, the Weierstrass $M$-Test is the main tool; on compact subintervals of the radius, it always applies.  Uniform convergence justifies swapping $\lim$ with $\int$ and $\frac{d}{dx}$ (with the extra hypothesis that the differentiated series converges uniformly).

\end{document}
